\documentclass[conference]{IEEEtran}
\IEEEoverridecommandlockouts
\usepackage[utf8]{inputenc}
\usepackage[spanish]{babel}
\usepackage{cite}
\usepackage{amsmath,amssymb,amsfonts}
\usepackage{algorithmic}
\usepackage{algorithm}
\usepackage{graphicx}
\usepackage{textcomp}
\usepackage{xcolor}
\usepackage{booktabs}
\usepackage{url}
\usepackage{hyperref}

\hypersetup{hidelinks}

\def\BibTeX{{\rm B\kern-.05em{\sc i\kern-.025em b}\kern-.08em
    T\kern-.1667em\lower.7ex\hbox{E}\kern-.125emX}}

% Macros de conveniencia
\newcommand{\FW}{Floyd-Warshall}
\newcommand{\INF}{\infty}

\begin{document}

\title{Paralelización del Algoritmo Floyd-Warshall mediante Ray:\\
Análisis de Rendimiento y Escalabilidad en Hardware de Alto Rendimiento}

\author{
\IEEEauthorblockN{Autor 1}
\IEEEauthorblockA{
  \textit{Instituto de Informática} \\
  \textit{Universidad Austral de Chile}\\
  Valdivia, Chile \\
  autor1@uach.cl
}
\and
\IEEEauthorblockN{Autor 2}
\IEEEauthorblockA{
  \textit{Instituto de Informática} \\
  \textit{Universidad Austral de Chile}\\
  Valdivia, Chile \\
  autor2@uach.cl
}
}

\maketitle

% ─────────────────────────────────────────────────────────────────────────────
\begin{abstract}
El algoritmo de \FW{} es uno de los algoritmos fundamentales de grafos
para el cálculo de caminos mínimos entre todos los pares de vértices, con
complejidad $O(n^3)$. Su estructura presenta restricciones de dependencia
que limitan la paralelización a la fase interna por pares $(i,j)$ para
cada iteración fija $k$. En este trabajo evaluamos experimentalmente el
impacto del framework Ray sobre la ejecución paralela de \FW{} en
hardware de alto rendimiento: AMD Ryzen Threadripper PRO 5975WX (32 núcleos)
con 128 GB de RAM. Diseñamos e implementamos dos versiones equivalentes del
algoritmo: una secuencial con operaciones vectorizadas NumPy y una paralela
utilizando actores Ray con particionamiento horizontal de la matriz. Ejecutamos
una batería de experimentos sistemáticos que cubren tamaños de entrada en
$n \in \{64, 128, 256, 512, 1024, 2048, 4096\}$, números de actores en
$\{1, 2, 4, 8, 16, 32\}$ y cinco grupos de escenarios que analizan speedup,
escalabilidad fuerte, escalabilidad débil, overhead de Ray y consumo de
recursos. Los resultados demuestran que [COMPLETAR CON DATOS REALES TRAS
EJECUTAR LOS BENCHMARKS].
\end{abstract}

\begin{IEEEkeywords}
Floyd-Warshall, Ray, computación paralela, HPC, paralelización por actores,
caminos mínimos, escalabilidad, benchmarking, Python
\end{IEEEkeywords}

% ─────────────────────────────────────────────────────────────────────────────
\section{Introducción}

El cálculo de caminos mínimos entre todos los pares de vértices (APSP, por
sus siglas en inglés) es una tarea central en teoría de grafos con aplicaciones
en redes de comunicaciones, sistemas de navegación, análisis de redes sociales
y optimización logística \cite{cormen2009introduction}. El algoritmo de \FW{}
\cite{floyd1962algorithm,warshall1962theorem} resuelve este problema de forma
exacta en tiempo $O(n^3)$ y espacio $O(n^2)$, siendo la solución de referencia
para grafos densos con pesos arbitrarios.

Con el auge de las arquitecturas multinúcleo y los sistemas distribuidos, la
paralelización de algoritmos clásicos mediante frameworks de alto nivel ha
ganado relevancia en la comunidad de computación de alto rendimiento (HPC).
Ray \cite{moritz2018ray} es un framework de ejecución distribuida para Python
que simplifica la programación paralela mediante el modelo de actores y la
gestión automática de un almacén de objetos compartidos (\textit{object store}).

La pregunta central de este trabajo es: \textit{¿en qué medida Ray permite
acelerar Floyd-Warshall, cuál es el speedup real, y cuál es el costo
computacional del paralelismo introducido?} Respondemos esta pregunta mediante
un diseño experimental riguroso que cuantifica no solo el tiempo de ejecución,
sino también la eficiencia paralela, el overhead de Ray, el consumo de memoria
y el uso de energía.

Las contribuciones principales de este artículo son:
\begin{itemize}
  \item Análisis formal de las dependencias de \FW{} y justificación de la
        estrategia de paralelización mediante actores Ray.
  \item Implementación reproducible de las versiones secuencial y paralela
        documentada con infraestructura Docker.
  \item Evaluación experimental con cinco grupos de escenarios sobre hardware
        de alto rendimiento (AMD Threadripper PRO 5975WX, 32 núcleos).
  \item Identificación cuantitativa del tamaño mínimo de problema a partir del
        cual Ray introduce beneficio neto.
\end{itemize}

El resto del artículo se organiza como sigue: la Sección~\ref{sec:relacionados}
revisa trabajos relacionados; la Sección~\ref{sec:solucion} presenta la solución
propuesta; la Sección~\ref{sec:metodologia} describe la metodología experimental;
la Sección~\ref{sec:escenarios} define los escenarios; la Sección~\ref{sec:resultados}
presenta los resultados; la Sección~\ref{sec:discusion} los discute; y la
Sección~\ref{sec:conclusiones} concluye.

% ─────────────────────────────────────────────────────────────────────────────
\section{Trabajos Relacionados}
\label{sec:relacionados}

La paralelización de \FW{} ha sido extensamente estudiada en arquitecturas
de memoria compartida y distribuida. Zhao et al.\ \cite{zhao2011parallel}
presentan una implementación en CUDA que aprovecha el paralelismo masivo de
la GPU para la fase interna, obteniendo speedups de hasta 40$\times$ sobre
la versión secuencial en CPU. En arquitecturas de CPU multinúcleo,
Venkataraman et al.\ \cite{venkataraman2003blocked} proponen la variante de
bloques (\textit{blocked Floyd-Warshall}) que mejora la localidad de caché
y facilita la paralelización con OpenMP.

En el contexto de sistemas distribuidos, Malewicz et al.\ \cite{malewicz2010pregel}
introducen Pregel para procesamiento de grafos a gran escala; sin embargo, su
modelo de mensajes por vértice no es eficiente para APSP de grafos densos.
Xia y Prasanna \cite{xia2004communication} analizan la complejidad de
comunicación del algoritmo en arquitecturas de paso de mensajes (MPI), demostrando
que el volumen de datos necesario por iteración $k$ es $O(n)$ si se mantienen
las filas distribuidas, lo que coincide con el diseño adoptado en este trabajo.

Ray \cite{moritz2018ray} fue evaluado originalmente en tareas de aprendizaje
reforzado; trabajos posteriores lo han utilizado para procesamiento de datos
a gran escala \cite{zaharia2012resilient} y simulaciones científicas. Sin
embargo, no encontramos estudios previos que evalúen sistemáticamente el uso
de actores Ray para \FW{} con análisis de overhead, escalabilidad y consumo
energético en hardware moderno.

% ─────────────────────────────────────────────────────────────────────────────
\section{Solución Propuesta}
\label{sec:solucion}

\subsection{Análisis de Dependencias del Algoritmo}

El algoritmo de \FW{} (Algoritmo~\ref{alg:fw}) itera sobre $n$ vértices
intermedios $k$. Para cada $k$, actualiza toda la matriz de distancias:

\begin{equation}
  d[i][j] \leftarrow \min(d[i][j],\; d[i][k] + d[k][j])
  \label{eq:fw}
\end{equation}

El bucle sobre $k$ es \textbf{secuencial obligatorio}: la iteración $k+1$
requiere los valores $d[i][k]$ y $d[k][j]$ actualizados en la iteración $k$.
Por el contrario, para un $k$ fijo, todas las actualizaciones $(i,j)$ son
independientes entre sí, ya que $d[i][k]$ y $d[k][j]$ son \textit{de solo
lectura} durante esa iteración \cite{cormen2009introduction}.

\begin{algorithm}
\caption{Floyd-Warshall}
\label{alg:fw}
\begin{algorithmic}
\FOR{$k = 0$ \TO $n-1$}
  \FOR{$i = 0$ \TO $n-1$} \hfill \textit{// paralelizable}
    \FOR{$j = 0$ \TO $n-1$}
      \STATE $d[i][j] \leftarrow \min(d[i][j], d[i][k] + d[k][j])$
    \ENDFOR
  \ENDFOR
\ENDFOR
\end{algorithmic}
\end{algorithm}

Esta estructura impone un límite teórico de speedup: el bucle externo (n
iteraciones) no puede paralelizarse, pero el trabajo interior ($O(n^2)$ por
iteración) sí puede distribuirse entre $p$ procesadores con un speedup
máximo de $O(p)$ para la parte paralelizable.

\subsection{Estrategia de Paralelización con Ray}

Adoptamos una \textbf{paralelización por actores con particionamiento
horizontal} de la matriz. Se crean $p$ actores Ray, cada uno manteniendo en
memoria local una partición de $\lceil n/p \rceil$ filas de la matriz global.

Para cada iteración $k$:
\begin{enumerate}
  \item El actor que posee la fila $k$ la publica en el \textit{object store}
        de Ray (vector de $n$ elementos, transferencia $O(n)$).
  \item Todos los actores reciben una referencia a la fila $k$; Ray provee
        acceso mediante memoria compartida sin copias adicionales.
  \item Cada actor actualiza su bloque de filas aplicando~\eqref{eq:fw}
        con operaciones vectorizadas NumPy.
  \item El bucle de $k$ continúa tras recibir la confirmación de todos los actores.
\end{enumerate}

Esta estrategia minimiza la comunicación a $O(n)$ datos por iteración,
comparado con $O(n^2/p)$ que implicaría transferir los bloques completos.
La columna $k$ es extraída localmente por cada actor desde su partición propia.

\subsection{Implementación}

La versión secuencial emplea vectorización con broadcasting de NumPy:
\texttt{dist = np.minimum(dist, col\_k[:, np.newaxis] + row\_k[np.newaxis, :])},
evitando bucles Python explícitos para maximizar el rendimiento de la línea
base. La versión paralela define la clase \texttt{FilasActor} anotada con
\texttt{@ray.remote}, que mantiene el bloque de filas como atributo de estado
y expone métodos \texttt{obtener\_fila}, \texttt{actualizar\_bloque} y
\texttt{obtener\_resultado}.

% ─────────────────────────────────────────────────────────────────────────────
\section{Metodología}
\label{sec:metodologia}

\subsection{Hardware y Software}

Los experimentos se ejecutan sobre un sistema con las siguientes
características: procesador AMD Ryzen Threadripper PRO 5975WX (32 núcleos
físicos, 64 hilos lógicos), 128 GB DDR4-3200, almacenamiento SSD NVMe 2~TB,
GPU NVIDIA TITAN V (12 GB HBM2), sistema operativo Linux (kernel 6.x).
El software utiliza Python 3.11, Ray 2.44, NumPy 1.26 y SciPy 1.12.

\subsection{Protocolo de Medición}

Todas las implementaciones se ejecutan dentro de un contenedor Docker
construido desde la imagen publicada en GHCR (Sección~\ref{sec:solucion}),
garantizando reproducibilidad del entorno. Para cada configuración se
realizan 10 repeticiones independientes con semillas distintas. Los tiempos
se miden con \texttt{time.perf\_counter()} que ofrece resolución
subnanosegundo en Linux.

Se aplica el test de Grubbs \cite{grubbs1969procedures} a nivel $\alpha=0.05$
para detectar y eliminar valores atípicos antes de calcular estadísticos.
Los intervalos de confianza al 95\% utilizan la distribución $t$ de Student
con $n-1$ grados de libertad.

\subsection{Métricas}

\begin{itemize}
  \item \textbf{Speedup}: $S(p) = T_{seq} / T_{ray}(p)$
  \item \textbf{Eficiencia paralela}: $E(p) = S(p) / p$
  \item \textbf{Overhead Ray}: $T_{overhead} = T_{setup} + T_{rearmado}$
  \item \textbf{Escalabilidad fuerte}: $S(p)$ con $n$ fijo, $p$ variable
  \item \textbf{Escalabilidad débil}: tiempo normalizado con $n \propto \sqrt{p}$
\end{itemize}

El uso de CPU se mide con \texttt{psutil.cpu\_percent(percpu=True)} cada
0.5 s. La energía CPU se obtiene desde
\texttt{/sys/class/powercap/} (interfaz RAPL del kernel Linux). El estado
de GPU se registra mediante \texttt{pynvml} (NVML).

% ─────────────────────────────────────────────────────────────────────────────
\section{Escenarios Experimentales}
\label{sec:escenarios}

Se definen cinco grupos de escenarios:

\begin{itemize}
  \item \textbf{E1 — Comparación directa}: Secuencial vs.\ Ray (32 actores)
        para $n \in \{64, 128, 256, 512, 1024, 2048, 4096\}$.
  \item \textbf{E2 — Escalabilidad fuerte}: $n$ fijo (2048), actores en
        $\{1, 2, 4, 8, 16, 32\}$.
  \item \textbf{E3 — Escalabilidad débil}: carga por actor constante,
        $n = 256 \cdot \sqrt{p}$.
  \item \textbf{E4 — Overhead}: matrices pequeñas ($n \leq 512$) para
        identificar el punto de crossover.
  \item \textbf{E5 — Carga máxima}: $n = 4096$ con todas las configuraciones.
\end{itemize}

El total de configuraciones únicas es de aproximadamente 120, con 10
repeticiones cada una (aprox.\ 1\,200 ejecuciones individuales).

% ─────────────────────────────────────────────────────────────────────────────
\section{Resultados}
\label{sec:resultados}

\textit{Nota: Las tablas y figuras de esta sección se generan automáticamente
a partir de los resultados del benchmark mediante los scripts del repositorio.
Ejecutar \texttt{make benchmark \&\& make analisis} para obtener los valores
reales.}

\subsection{Comparación de Tiempos (E1)}

La Tabla~\ref{tab:tiempos} presenta los tiempos de ejecución medios con
intervalos de confianza al 95\% para ambas implementaciones.

\begin{table}[htbp]
\caption{Tiempos de ejecución: Secuencial vs.\ Ray (32 actores)}
\label{tab:tiempos}
\begin{center}
\begin{tabular}{|r|c|c|}
\hline
\textbf{$n$} & \textbf{$T_{seq}$ (s)} & \textbf{$T_{ray}$ (s)} \\
\hline
% Los valores se completan automáticamente con: make tablas
% Resultado en: resultados/tabla_tiempos.tex
\input{resultados/tabla_tiempos.tex}
\end{tabular}
\end{center}
\end{table}

\subsection{Speedup y Eficiencia Paralela (E1, E2)}

La Tabla~\ref{tab:speedup} resume el speedup $S(p)$ y la eficiencia
paralela $E(p)$ para las configuraciones más representativas.

\begin{table}[htbp]
\caption{Speedup y eficiencia paralela}
\label{tab:speedup}
\begin{center}
\begin{tabular}{|r|r|c|c|r|r|}
\hline
\textbf{$n$} & \textbf{$p$} & \textbf{$T_{seq}$ (s)} & \textbf{$T_{ray}$ (s)} & \textbf{$S(p)$} & \textbf{$E(p)$} \\
\hline
\input{resultados/tabla_speedup.tex}
\end{tabular}
\end{center}
\end{table}

\subsection{Overhead de Ray (E4)}

La Tabla~\ref{tab:overhead} descompone el tiempo de ejecución Ray en tiempo
de cómputo efectivo y overhead de coordinación.

\begin{table}[htbp]
\caption{Descomposición del overhead de Ray}
\label{tab:overhead}
\begin{center}
\begin{tabular}{|r|r|c|c|r|}
\hline
\textbf{$n$} & \textbf{$p$} & \textbf{$T_{cómp.}$} & \textbf{$T_{overhead}$} & \textbf{\%} \\
\hline
\input{resultados/tabla_overhead.tex}
\end{tabular}
\end{center}
\end{table}

\subsection{Escalabilidad Fuerte y Débil (E2, E3)}

\begin{figure}[htbp]
\centerline{\includegraphics[width=\columnwidth]{graficos/speedup_vs_workers.pdf}}
\caption{Escalabilidad fuerte: speedup vs.\ número de actores Ray para
distintos tamaños de matriz. La línea discontinua representa el speedup ideal.}
\label{fig:speedup_workers}
\end{figure}

\begin{figure}[htbp]
\centerline{\includegraphics[width=\columnwidth]{graficos/speedup_vs_tamano.pdf}}
\caption{Speedup vs.\ tamaño de matriz $n$ para distintas cantidades de
actores Ray. Se observa el punto de crossover a partir del cual Ray supera
al secuencial.}
\label{fig:speedup_tamano}
\end{figure}

\begin{figure}[htbp]
\centerline{\includegraphics[width=\columnwidth]{graficos/escalabilidad_debil.pdf}}
\caption{Escalabilidad débil: tiempo de ejecución normalizado respecto al
caso de un solo actor, manteniendo carga por actor constante.}
\label{fig:escal_debil}
\end{figure}

\subsection{Consumo de Recursos (E1, E5)}

\begin{figure}[htbp]
\centerline{\includegraphics[width=\columnwidth]{graficos/consumo_cpu_ram.pdf}}
\caption{Uso de CPU y RAM pico para versión secuencial vs.\ Ray (32 actores)
por tamaño de matriz.}
\label{fig:recursos}
\end{figure}

La Tabla~\ref{tab:recursos} detalla el uso promedio de CPU, RAM pico y
consumo energético estimado por configuración.

\begin{table}[htbp]
\caption{Métricas de recursos por configuración}
\label{tab:recursos}
\begin{center}
\begin{tabular}{|r|c|r|r|r|r|}
\hline
\textbf{$n$} & \textbf{Alg.} & \textbf{$p$} & \textbf{CPU\%} & \textbf{RAM (MB)} & \textbf{E (J)} \\
\hline
\input{resultados/tabla_recursos.tex}
\end{tabular}
\end{center}
\end{table}

% ─────────────────────────────────────────────────────────────────────────────
\section{Discusión}
\label{sec:discusion}

\subsection{Speedup y Eficiencia}

[COMPLETAR CON ANÁLISIS DE LOS RESULTADOS REALES]

Los resultados de la Tabla~\ref{tab:speedup} muestran que el speedup de
Ray es significativo para matrices grandes ($n \geq$ [VALOR]). Para matrices
pequeñas ($n < $ [VALOR]), el overhead de Ray ($T_{overhead} \approx$ [VALOR]~s)
supera al beneficio del paralelismo, lo que se confirma en la Tabla~\ref{tab:overhead}.

\subsection{Escalabilidad}

La Fig.~\ref{fig:speedup_workers} evidencia [COMPLETAR]. La Ley de Amdahl
predice un speedup máximo de $1/(1-\beta)$ donde $\beta$ es la fracción
paralelizable. En Floyd-Warshall, la fracción paralelizable por iteración es
cercana a 1 para la parte $(i,j)$, pero las $n$ iteraciones de $k$ requieren
sincronización secuencial.

\subsection{Overhead de Coordinación}

Ray introduce tres fuentes de overhead: (i) serialización y transferencia de
la fila $k$ al \textit{object store}, (ii) scheduling de tareas en los actores,
y (iii) sincronización mediante \texttt{ray.get}. El overhead absoluto es
aproximadamente constante por iteración y representa un porcentaje decreciente
del tiempo total para $n$ grande.

\subsection{Amenazas a la Validez}

\textit{Validez interna}: Los tiempos pueden verse afectados por procesos del
sistema operativo. Se mitiga con 10 repeticiones y el test de Grubbs.
\textit{Validez externa}: Los resultados son específicos al hardware AMD
Threadripper. La generalización a otros sistemas requiere replicación.
\textit{Validez de constructo}: El speedup mide el tiempo de reloj completo,
incluyendo overhead de Ray. Esto refleja el beneficio neto para el usuario final.

% ─────────────────────────────────────────────────────────────────────────────
\section{Conclusiones}
\label{sec:conclusiones}

En este trabajo evaluamos experimentalmente la paralelización del algoritmo
de Floyd-Warshall mediante actores Ray sobre hardware de alto rendimiento
(AMD Threadripper PRO 5975WX, 32 núcleos). Las principales conclusiones son:

\begin{itemize}
  \item [COMPLETAR TRAS EJECUTAR BENCHMARKS]
  \item El overhead de Ray es [VALOR]~s y permanece aproximadamente constante
        independientemente del tamaño de la matriz.
  \item Ray ofrece beneficio neto para $n \geq$ [VALOR], con un speedup de
        hasta [VALOR]$\times$ con 32 actores.
  \item La eficiencia paralela decrece con el número de actores, lo que
        es consistente con la ley de Amdahl aplicada a las $n$ barreras de
        sincronización del algoritmo.
  \item El consumo de memoria RAM escala linealmente con el número de actores
        debido a la replicación de la fila $k$ en el \textit{object store}.
\end{itemize}

Como trabajo futuro se propone: (i) implementar la variante de bloques
(\textit{blocked Floyd-Warshall}) para mejor localidad de caché; (ii)
extender la evaluación a múltiples nodos con Ray en modo distribuido real;
(iii) comparar con implementaciones OpenMP/MPI equivalentes.

El entorno experimental completo está disponible públicamente en:\\
\url{https://github.com/martinmaza/floyd-warshall-ray}\\
con imagen Docker reproducible en:\\
\url{ghcr.io/martinmaza/floyd-warshall-ray:latest}

% ─────────────────────────────────────────────────────────────────────────────
\section*{Agradecimientos}

Los autores agradecen al Instituto de Informática de la Universidad Austral
de Chile por el acceso al hardware de alto rendimiento utilizado en este estudio.

% ─────────────────────────────────────────────────────────────────────────────
\begin{thebibliography}{00}

\bibitem{floyd1962algorithm}
R.~W. Floyd, ``Algorithm 97: Shortest path,''
\textit{Commun. ACM}, vol.~5, no.~6, p.~345, Jun. 1962.

\bibitem{warshall1962theorem}
S.~Warshall, ``A theorem on boolean matrices,''
\textit{J. ACM}, vol.~9, no.~1, pp.~11--12, Jan. 1962.

\bibitem{cormen2009introduction}
T.~H. Cormen, C.~E. Leiserson, R.~L. Rivest, and C.~Stein,
\textit{Introduction to Algorithms}, 3rd~ed.
Cambridge, MA: MIT Press, 2009.

\bibitem{moritz2018ray}
P.~Moritz \textit{et al.}, ``Ray: A distributed framework for emerging AI
applications,'' in \textit{Proc. 13th USENIX Symp. Operating Systems Design
and Implementation (OSDI)}, 2018, pp.~561--577.

\bibitem{zhao2011parallel}
Y.~Zhao and P.~Tao, ``A parallel Floyd-Warshall algorithm on a GPU,''
in \textit{Proc. Int. Conf. Parallel and Distributed Computing, Applications
and Technologies (PDCAT)}, 2011, pp.~299--305.

\bibitem{venkataraman2003blocked}
G.~Venkataraman, S.~Sahni, and S.~Mukhopadhyaya, ``A blocked all-pairs
shortest-paths algorithm,'' \textit{J. Exp. Algorithmics}, vol.~8, 2003.

\bibitem{malewicz2010pregel}
G.~Malewicz \textit{et al.}, ``Pregel: A system for large-scale graph
processing,'' in \textit{Proc. ACM SIGMOD Int. Conf. Management of Data},
2010, pp.~135--146.

\bibitem{xia2004communication}
Y.~Xia and V.~K. Prasanna, ``Communication optimal parallel shortest path
algorithms,'' in \textit{Proc. Int. Parallel and Distributed Processing Symp.
(IPDPS)}, 2004.

\bibitem{zaharia2012resilient}
M.~Zaharia \textit{et al.}, ``Resilient distributed datasets: A fault-tolerant
abstraction for in-memory cluster computing,'' in \textit{Proc. 9th USENIX
Conf. Networked Systems Design and Implementation (NSDI)}, 2012.

\bibitem{grubbs1969procedures}
F.~E. Grubbs, ``Procedures for detecting outlying observations in samples,''
\textit{Technometrics}, vol.~11, no.~1, pp.~1--21, 1969.

\end{thebibliography}

\end{document}
